% =============================================================================
%  lec27.tex —— 第 27 讲 E_∞-环（底本 PDF 第 199–203 页）
% =============================================================================
\lecture{27}{$E_\infty$-环}
\begin{center}\small\color{black!55}（底本 PDF 第 199–203 页）\end{center}

本讲解释下述类比：
\begin{center}
\begin{tabular}{ll}
\hline
经典代数 & $\infty$-范畴代数 \\
\hline
集合 & 空间 \\
阿贝尔群 & 谱 \\
张量积 & smash 积 \\
结合环 & $E_{1}$-环 \\
交换环 & $E_{\infty}$-环 \\
\hline
\end{tabular}
\end{center}

\yizhu{本讲的“$E_\infty$-环”即 HA 中的“交换环谱”（commutative ring spectrum）；采用 $E_\infty$ 的说法，是为强调其结构由 $E_\infty$-操作子的代数给出，而不是在 $1$-范畴层面要求严格交换。对照表最后一行把“交换环”对到 $E_\infty$-环，正是这个意思。\textbf{[来源：A]}}

\section{谱与阿贝尔群}

\nitem\label{L27-1.1}
本节解释谱与阿贝尔群之间的类比。

\nitem\label{L27-1.2}
阿贝尔群是 $\Set$ 中的交换幺半群 $G$，且每个元素都有逆。后一条件说的是映射
\[
(p_{1},m)\colon G\times G\to G\times G,\qquad (m,p_{2})\colon G\times G\to G\times G
\]
可逆。此概念在任意笛卡尔对称 monoidal $\infty$-范畴中都有意义。

\begin{definition}\label{L27-1.3}
设 $\cC$ 是承认有限乘积的 $\infty$-范畴。若结合幺半群 $G\in\Mon_{E_1}(\cC)$ 的态射 $(p_{1},m)$ 与 $(m,p_{2})$ 可逆，则称 $G$ 为 $\cC$ 中的\textbf{群}。

对 $1\le k\le\infty$，若 $E_{k}$-幺半群 $G\in\Mon_{E_k}(\cC)$ 在
\[
\Mon_{E_k}(\cC)\to\Mon_{E_1}(\cC)
\]
下的像是一个群，则称 $G$ 为 $\cC$ 中的 \textbf{$E_{k}$-群}\footnote{另一种说法是“群样的 $E_k$-幺半群”（grouplike $E_k$-monoids）。}。令
\[
\Grp_{E_k}(\cC)\subseteq\Mon_{E_k}(\cC)
\]
为 $\cC$ 中 $E_{k}$-群构成的全子 $\infty$-范畴。
\end{definition}

\begin{example}\label{L27-1.4}
设 $X\in\Spc_{*}$ 是任意带基点的空间。对 $0<k<\infty$，$k$ 重环路空间 $\Omega^{k}X$ 带有如下的典范 $E_{k}$-群结构。

选取 $X$ 的一个拓扑实现，则 $\Omega^{k}X$ 可以实现为满足“$f$ 把边界送到 $X$ 的基点”的映射 $f\colon[-1,1]^{k}\to X$ 构成的拓扑空间。对 ${}^{t}E_{k}$ 中的任意 $m$ 元运算 $\gamma\colon\square^{k}\times\langle m\rangle^{\circ}\to\square^{k}$，有同态
\[
\mult_{\gamma}\colon(\Omega^{k}X)^{\times m}\to\Omega^{k}X,
\]
它把 $(f_{1},\dots,f_{m})\in(\Omega^{k}X)^{\times m}$ 送到映射
\[
f(s)\defeq
\begin{cases}
f_{i}(t), & s=\gamma(t,i),\\
*, & \text{其他情形}.
\end{cases}
\]
\end{example}

可以验证，这在 $\Top$ 中定义了 ${}^{t}E_{k}$-群 $\Omega^{k}X$（其中一切都在拓扑空间上充实）。过渡到 $\infty$-范畴，便得到 $\Spc$ 中的 $E_{k}$-群。上述构造对 $X$ 具有函子性，并给出函子
\begin{equation}\label{L27-1.1eq}
\Omega^{k}\colon\Spc_{*}\to\Grp_{E_k}(\Spc).
\end{equation}

\begin{exercise}\label{L27-1.5}
对任意承认有限乘积的 $\infty$-范畴 $\cC$，构造函子 $\Omega^{k}\colon\cC_{*}\to\Grp_{E_k}(\cC)$。提示：Yoneda 引理。
\end{exercise}

\nitem\label{L27-1.6}
函子 \eqref{L27-1.1eq} 不是等价，甚至不是保守的。事实上，$\Omega X$ 只能看见 $X$ 的中性连通分支。结果表明，这正是 \eqref{L27-1.1eq} 不是等价的唯一原因。

\begin{theorem}[Boardman--Vogt--May]\label{L27-1.7}
对 $0<k<\infty$，令 $(\Spc_{*})_{\ge k}\subseteq\Spc_{*}$ 为 $k$-连通空间\footnote{即对 $i<k$ 有 $\pi_{i}\simeq 0$。}构成的全子 $\infty$-范畴。则函子
\begin{equation}\label{L27-1.2eq}
\Omega^{k}\colon(\Spc_{*})_{\ge k}\to\Grp_{E_k}(\Spc)
\end{equation}
是等价。
\end{theorem}

\begin{remark}\label{L27-1.8}
退化情形 $k=0$ 是：$\Spc_{*}\simeq\Mon_{E_0}(\Spc)$。
\end{remark}

\begin{exercise}\label{L27-1.9}
对 $0<k<\infty$，构造典范交换图
\[
\begin{tikzcd}[column sep=4.2em,row sep=3.2em]
(\Spc_{*})_{\ge k+1} \arrow[r,"\Omega"] \arrow[d,"\Omega^{k+1}"] & (\Spc_{*})_{\ge k} \arrow[d,"\Omega^{k}"] \\
\Grp_{E_{k+1}}(\Spc) \arrow[r] & \Grp_{E_k}(\Spc).
\end{tikzcd}
\]
\end{exercise}

\nitem\label{L27-1.10}
由上述练习，可以对 $k$ 取逆向极限，得到连通谱\footnote{即对 $i<0$ 有 $\pi_{i}\simeq 0$。}的下述刻画。

\begin{corollary}\label{L27-1.11}
存在典范等价
\[
\Omega^{\infty}\colon\Sptr_{\ge 0}\xrightarrow{\ \simeq\ }\Grp_{E_\infty}(\Spc),
\]
它与到 $\Spc$ 的遗忘函子相容。
\end{corollary}

\begin{remark}\label{L27-1.12}
换言之，连通谱正是 Picard $\infty$-群胚，即每个对象关于张量积都可逆的对称 monoidal $\infty$-群胚。
\end{remark}

\begin{remark}\label{L27-1.13}
把 \eqref{L27-1.2eq} 的逆记为
\[
B^{k}\colon\Grp_{E_k}(\Spc)\to(\Spc_{*})_{\ge k},
\]
并称之为 \textbf{$k$ 重解环函子}。粗略地说，此函子可如下描述。

对 $k=1$，给定笛卡尔 monoidal $\infty$-范畴 $\cC$ 中的任意结合幺半群 $G$，可以作 Bar 构造
\[
*\times *\rightrightarrows *\times G\times *\rightrightarrows *\times G\times G\times *\cdots,
\]
它是 $\cC$ 中的单纯对象。该图的余极限 $BG$（若存在）计算 $*$ 作为 $G$-双模与自身的相对张量积。可以验证，所得函子
\[
B\colon\Grp_{E_1}(\Spc)\to\Spc_{*}
\]
的像包含在 $(\Spc_{*})_{\ge 0}$ 中，并且确实给出所需的解环函子。

对一般的 $0<k<\infty$，注意复合
\[
\Grp_{E_1}(\Spc)\xrightarrow{\ B\ }(\Spc_{*})_{\ge 1}\to\Spc
\]
与有限乘积交换，因此诱导函子
\[
B\colon\Grp_{E_{k+1}}(\Spc)\simeq\Mon_{E_k}(\Grp_{E_1}(\Spc))\to\Mon_{E_k}(\Spc),
\]
其中第一个等价来自 Dunn 可加性定理。可以验证此函子分解为
\[
B\colon\Grp_{E_{k+1}}(\Spc)\to\Grp_{E_k}(\Spc).
\]
于是得到函子
\[
B^{k}\colon\Grp_{E_k}(\Spc)\to\Spc_{*}.
\]
可以验证其像包含在 $(\Spc_{*})_{\ge k}$ 中，并且确实给出所需的 $k$ 重解环函子。
\end{remark}

\section{smash 积与张量积}

\nitem\label{L27-2.1}
本节解释谱的 smash 积与阿贝尔群的张量积之间的类比。

\nitem\label{L27-2.2}
回忆有函子
\[
\pi_{\bullet}\colon\Sptr\to\mathrm{grAb},
\]
它把谱 $X$ 送到其同伦群
\[
\pi_{m}X\simeq\pi_{0}\Map_{\Sptr}(S[m],X).
\]

\begin{construction}\label{L27-2.3}
考虑 $\Sptr$ 中由单位结构给出的同构 $S\otimes S\xrightarrow{\ \simeq\ }S$。利用 $-\otimes-$ 与每个因子的余极限交换这一事实，对整数 $m,n$ 得到同构
\[
S[m]\otimes S[n]\xrightarrow{\ \simeq\ }S[m+n].
\]
因此对谱 $X,Y$，得到空间之间的态射
\[
\Map_{\Sptr}(S[m],X)\times\Map_{\Sptr}(S[n],Y)\to\Map_{\Sptr}(S[m]\otimes S[n],X\otimes Y)
\xrightarrow{\ \simeq\ }\Map_{\Sptr}(S[m+n],X\otimes Y).
\]
取 $\pi_{0}$，得到映射
\[
\pi_{m}X\times\pi_{n}Y\to\pi_{m+n}(X\otimes Y).
\]
可以验证此映射双线性，因此给出 $\Ab$ 中的态射
\[
\pi_{m}X\otimes\pi_{n}Y\to\pi_{m+n}(X\otimes Y),
\]
它们可以组装为态射
\[
\pi_{\bullet}X\otimes\pi_{\bullet}Y\to\pi_{\bullet}(X\otimes Y),
\]
其中源为分次张量积。可以验证，这给出函子 $\pi_{\bullet}$ 上的右 lax 对称 monoidal 结构\footnote{此结构中没有更高阶数据，因为 $\mathrm{grAb}$ 是通常范畴。}。换言之，有 $\infty$-操作子映射
\[
\pi_{\bullet}\colon\Sptr^{\otimes}\to\mathrm{grAb}^{\otimes}.
\]
\end{construction}

\begin{warning}\label{L27-2.4}
函子 $\pi_{\bullet}$ 不是对称 monoidal 的。事实上，$\pi_{\bullet}(S)$ 与 $\Z$ 大不相同。
\end{warning}

\nitem\label{L27-2.5}
下述练习说明，限制到 $\Sptr_{\ge 0}$ 时 $\pi_{0}$ 是对称 monoidal 的。

\begin{exercise}\label{L27-2.6}
全子 $\infty$-范畴 $\Sptr_{\ge 0}\subseteq\Sptr$ 对 smash 积封闭。提示：$\Sptr_{\ge 0}$ 由 $S$ 在余极限与扩张下生成。
\end{exercise}

\begin{exercise}\label{L27-2.7}
$\pi_{0}\colon\Sptr_{\ge 0}\to\Ab$ 上的 lax 对称 monoidal 结构是严格的。
\end{exercise}

\begin{exercise}\label{L27-2.8}
用相应的 Picard $\infty$-群胚描述谱的 smash 积。提示：可以先看 $E_{0}$ 情形。
\end{exercise}

\section{$E_{k}$-环}

\begin{definition}\label{L27-3.1}
给 $\Sptr$ 装备 smash 积对称 monoidal 结构。对 $0\le k\le\infty$，\textbf{$E_{k}$-环}是 $\Sptr$ 中的 $E_{k}$-代数。
\end{definition}

\begin{construction}\label{L27-3.2}
由于 $\Sptr\to\mathrm{grAb}$ 自然右 lax 对称 monoidal，得到函子
\[
\Alg_{E_k}(\Sptr)\to\Alg_{E_k}(\mathrm{grAb}).
\]
因此：
\begin{itemize}[leftmargin=2.2em,itemsep=0.1em]
\item 对 $E_{1}$-环 $A$，得到结合环 $\pi_{\bullet}A$；
\item 对 $k\ge 2$ 的 $E_{k}$-环 $A$，得到分次交换环 $\pi_{\bullet}A$。
\end{itemize}
\end{construction}

\begin{construction}\label{L27-3.3}
注意 $\pi_{0}\colon\Sptr_{\ge 0}\to\Ab$ 可以通过 $\Ab\simeq\Sptr^{\heartsuit}$ 等同于函子 $\tau_{\le 0}$，并且它有右伴随
\[
\Ab\to\Sptr_{\ge 0},\qquad R\mapsto HR,
\]
把 $R$ 送到 Eilenberg--MacLane 谱 $HR$。由于左伴随函子自然对称 monoidal，右伴随带有右 lax 对称 monoidal 结构。特别地，得到函子
\[
\Alg_{E_k}(\Ab)\to\Alg_{E_k}(\Sptr_{\ge 0})\to\Alg_{E_k}(\Sptr).
\]
换言之：
\begin{itemize}[leftmargin=2.2em,itemsep=0.1em]
\item 对结合环 $R$，得到 $HR$ 上的 $E_{1}$-环结构；
\item 对交换环 $R$，得到 $HR$ 上的 $E_{\infty}$-环结构。
\end{itemize}
\end{construction}

\nitem\label{L27-3.4}
最后一项任务是把 $HR$-模与 $R$-模联系起来。为简便起见，我们在 $E_{\infty}$ 情形下工作。

从现在起，设 $R$ 是交换环，$\Mod_{R}$ 是 $R$-模范畴（阿贝尔范畴），$\Mod_{HR}(\Sptr)$ 是 $\Sptr$ 中 $HR$-模的 $\infty$-范畴。

\begin{construction}\label{L27-3.5}
考虑 $\Mod_{R}$ 上的相对张量积对称 monoidal 结构。过渡到导出 $\infty$-范畴，得到 $D(R)\defeq D(\Mod_{R})$ 上的可表现对称 monoidal 结构（细节见 HA.7.1.2.12）。由于 $D(R)$ 稳定，得到交换代数对象 $D(R)\in\Alg_{\mathrm{Comm}}(\PrSt^{L})$。回忆 $\Sptr$ 是 $\PrSt^{L}$ 的单位。于是得到 $\Alg_{\mathrm{Comm}}(\PrSt^{L})$ 中的态射 $\Sptr\to D(R)$。换言之，一个对称 monoidal 函子
\[
F\colon\Sptr\to D(R).
\]
\end{construction}

\nitem\label{L27-3.6}
注意底函子 $\Sptr\to D(R)$ 是唯一的保余极限函子，它把 $S\in\Sptr$ 送到 $R\in D(R)$。另一方面，复合
\[
G\colon D(R)\to D(\Z)\xrightarrow{\ DK\ }\Sptr
\]
的左伴随也把 $S$ 送到 $R$。于是得到伴随
\[
F\colon\Sptr\rightleftarrows D(R)\colon G,
\]
其中 $F$ 对称 monoidal，而 $G$ 右 lax 对称 monoidal\footnote{更确切地说，我们有在 $\mathrm{Fin}_{*}$ 上定义的伴随 $\Sptr^{\otimes}\rightleftarrows D(R)^{\otimes}$。}。注意 $G$ 也保持余极限。

\begin{exercise}\label{L27-3.7}
证明复合
\[
G\circ F\colon\Sptr\to D(R)\to\Sptr
\]
等同于 $HR\otimes-$。
\end{exercise}

\begin{theorem}[HA.7.1.2.13]\label{L27-3.8}
存在对称 monoidal $\infty$-范畴的典范等价
\[
\Mod_{HR}(\Sptr)\simeq D(R),
\]
使得伴随 $F\colon\Sptr\rightleftarrows D(R)\colon G$ 可以等同于沿单位态射 $u\colon S\to HR$ 的
\[
\ind_{u}\colon\Mod_{S}(\Sptr)\rightleftarrows\Mod_{HR}(\Sptr)\colon\res_{u}.
\]
\end{theorem}

\appendixsection{Barr--Beck--Lurie 定理}

\yizhu{“$G$-分裂单纯余极限”指：对单纯对象 $V\in\cD_{\Delta}$，若 $G(V)$ 在 $\cC$ 中带有分裂（即其面/退化结构额外给出收缩），则 $V$ 在 $\cD$ 中有余极限，且 $G$ 保持该余极限。这是 HA.4.7.5.3 所用“Beck 条件”的 $\infty$-版本，替代了经典 Barr--Beck 定理中“$G$ 反射等化子”的要求。\textbf{[来源：A]}}

\nitem\label{L27-A.1}
定理 \ref{L27-3.8} 的证明的关键输入是下述定理。

\begin{theorem}[Barr--Beck--Lurie]\label{L27-A.2}
设 $F\colon\cC\rightleftarrows\cD\colon G$ 是 $\infty$-范畴之间的伴随，且满足
\begin{enumerate}[label=(\roman*),leftmargin=2.4em,itemsep=0.1em]
\item 函子 $G$ 保守；
\item 函子 $G$ 保持 $G$-分裂单纯余极限\footnote{这指的是：设 $V\in\cD_{\Delta}\defeq\Fun(\Delta^{\op},\cD)$ 是使得 $G(V)$ 承认分裂的单纯对象。则 $V$ 在 $\cD$ 中承认余极限，且该余极限被 $G$ 保持。}。
\end{enumerate}
则存在典范等价 $\cD\simeq\mathrm{LMod}_{T}(\cC)$，其中 $T$ 是作用在 $\cC$ 上的单子，其底自同态为 $G\circ F$。此外，伴随 $F\colon\cC\rightleftarrows\cD\colon G$ 可以等同于
\[
\ind_{T}\colon\cC\to\mathrm{LMod}_{T}(\cC)\colon\oblv_{T}.
\]
\end{theorem}

\nitem\label{L27-A.3}
\textbf{建议阅读}\quad HA.4.7。
